% preamble-exam-zh.tex — 共享 preamble
% 用于所有 exam-zh 模板
%
% 样式策略：
%   屏幕 → 语义彩色（蓝=关键想法，红=易错，绿=路线，灰=提示/教师备注）
%   打印 → 自动降级为黑白（通过 print mode 或 monochrome 选项）

\documentclass{exam-zh}

% ---------- 基础配置 ----------
\examsetup{
  page/size = a4paper,
  page/show-head = false,
  page/show-foot = false,
  sealline/show = false,
  font = times,
  math-font = xits,
}

\usepackage{tcolorbox}
\usepackage{needspace}
\tcbuselibrary{skins,breakable}

% ---------- 打印/屏幕颜色切换 ----------
% 用法：\documentclass[monochrome]{exam-zh} 即可切换为纯黑白
% 注意：不要使用 \if@xxx 放在 makeatother 外部；这里改成普通条件名。
\newif\ifeduprintcolor
\eduprintcolortrue

\makeatletter
\@ifclasswith{exam-zh}{monochrome}{\eduprintcolorfalse}{}
\makeatother

% ---------- 语义颜色定义 ----------
% 屏幕：有色彩区分；打印：降级为灰度
\ifeduprintcolor
  % --- 屏幕彩色模式 ---
  \definecolor{edu-blue}{RGB}{47, 82, 122}       % 关键想法
  \definecolor{edu-blue-bg}{RGB}{243, 246, 252}
  \definecolor{edu-red}{RGB}{180, 40, 40}         % 易错提醒
  \definecolor{edu-red-bg}{RGB}{255, 245, 240}
  \definecolor{edu-green}{RGB}{34, 100, 60}       % 路线图
  \definecolor{edu-green-bg}{RGB}{242, 250, 245}
  \definecolor{edu-orange}{RGB}{160, 90, 0}       % 训练目标
  \definecolor{edu-orange-bg}{RGB}{255, 248, 235}
  \definecolor{edu-gray}{RGB}{100, 100, 100}      % 提示/教师备注
  \definecolor{edu-gray-bg}{RGB}{248, 248, 248}
  \definecolor{edu-purple}{RGB}{90, 50, 120}      % 思考题
  \definecolor{edu-purple-bg}{RGB}{248, 244, 252}
\else
  % --- 打印黑白模式 ---
  \definecolor{edu-blue}{gray}{0.15}
  \definecolor{edu-blue-bg}{gray}{0.96}
  \definecolor{edu-red}{gray}{0.25}
  \definecolor{edu-red-bg}{gray}{0.97}
  \definecolor{edu-green}{gray}{0.20}
  \definecolor{edu-green-bg}{gray}{0.97}
  \definecolor{edu-orange}{gray}{0.25}
  \definecolor{edu-orange-bg}{gray}{0.98}
  \definecolor{edu-gray}{gray}{0.40}
  \definecolor{edu-gray-bg}{gray}{0.97}
  \definecolor{edu-purple}{gray}{0.30}
  \definecolor{edu-purple-bg}{gray}{0.98}
\fi

% ---------- 自定义教学环境 ----------
% 共性：enhanced + breakable（跨页不断裂）

% 原题卡片 — 强边框，突出显示
\newtcolorbox{problemcard}{
  enhanced, breakable,
  colback=edu-blue-bg,
  colframe=edu-blue,
  boxrule=1.2pt,
  arc=0pt,
  left=3mm, right=3mm, top=3mm, bottom=3mm,
}

% 解题路线图 — 左边线 + 浅底
\newtcolorbox{routebox}{
  enhanced, breakable,
  colback=edu-green-bg,
  colframe=edu-green!30,
  boxrule=0pt,
  borderline west={2.5pt}{0pt}{edu-green},
  left=3mm, right=3mm, top=3mm, bottom=3mm,
}

% 关键想法 — 蓝色左边线
\newtcolorbox{keyideabox}{
  enhanced, breakable,
  colback=edu-blue-bg,
  colframe=edu-blue!30,
  boxrule=0pt,
  borderline west={2.5pt}{0pt}{edu-blue},
  left=3mm, right=3mm, top=3mm, bottom=3mm,
}

% 易错提醒 — 红色左边线
\newtcolorbox{mistakebox}{
  enhanced, breakable,
  colback=edu-red-bg,
  colframe=edu-red!20,
  boxrule=0pt,
  borderline west={3pt}{0pt}{edu-red},
  left=3mm, right=3mm, top=2.5mm, bottom=2.5mm,
}

% 提示 — 灰色虚线左边线
\newtcolorbox{hintbox}{
  enhanced, breakable,
  colback=edu-gray-bg,
  colframe=edu-gray!20,
  boxrule=0pt,
  borderline west={2pt}{0pt}{edu-gray, dashed},
  left=3mm, right=3mm, top=2mm, bottom=2mm,
}

% 思考题 — 紫色虚线左边线
\newtcolorbox{thinkbox}{
  enhanced, breakable,
  colback=edu-purple-bg,
  colframe=edu-purple!20,
  boxrule=0pt,
  borderline west={2pt}{0pt}{edu-purple, dashed},
  left=2.5mm, right=2.5mm, top=2.5mm, bottom=2.5mm,
}

% 教师备注 — 灰色左边线，小字
\newtcolorbox{teachernote}{
  enhanced, breakable,
  colback=edu-gray-bg,
  colframe=edu-gray!15,
  boxrule=0pt,
  borderline west={2pt}{0pt}{edu-gray!60},
  left=2.5mm, right=2.5mm, top=2.5mm, bottom=2.5mm,
  fontupper=\small,
  colupper=black!60,
}

% 训练目标 — 橙色左边线，小字
\newtcolorbox{traininggoal}{
  enhanced, breakable,
  colback=edu-orange-bg,
  colframe=edu-orange!15,
  boxrule=0pt,
  borderline west={1.5pt}{0pt}{edu-orange},
  left=2mm, right=2mm, top=1mm, bottom=1mm,
  fontupper=\small,
}

% 答题区 — 无边框，纯留白
\newcommand{\answerarea}[1][40mm]{%
  \vspace{2mm}%
  \begin{tcolorbox}[
    enhanced, breakable,
    colback=white,
    colframe=white,
    boxrule=0pt,
    height=#1,
    valign=top,
    bottom=0mm,
    after skip=0mm,
  ]
  \end{tcolorbox}%
}

% ---------- 字体设置 ----------
% exam-zh (ctexbook) already sets CJK fonts via ctex.
% Only override if specific fonts are needed.
% macOS: Songti SC, STHeiti, STFangsong
% Linux: SimSun, SimHei, FangSong

% ---------- 页面设置 ----------
% exam-zh (ctexbook) already loads geometry.
% Override margins if needed:
% \geometry{top=16mm, bottom=16mm, left=14mm, right=14mm}

\examsetup{
  question/show-answer = true,
  fillin/show-answer = true,
  solution/show-solution = show-stay,
}

% ---------- 教师版解答样式 ----------
\newtcolorbox{solutionblock}{
  enhanced, breakable,
  colback=edu-blue-bg,
  colframe=edu-blue!25,
  boxrule=0pt,
  borderline west={2pt}{0pt}{edu-blue!50},
  arc=0pt,
  left=3mm, right=3mm, top=2mm, bottom=2mm,
  before skip=2mm,
  after skip=3mm,
  fontupper=\small,
}

% 解答步骤编号
\newcounter{solstep}
\newcommand{\solstep}[2]{%
  \stepcounter{solstep}%
  \par\medskip
  \noindent\hangindent=6mm
  {\color{edu-blue}\bfseries\textcircled{\scriptsize\thesolstep}\enspace #1}\enspace #2\par
}

\begin{document}

\section*{求一次函数解析式专题（二） · 练习卷（教师版）}

\section*{一、选择题（每题 12 分，共 48 分）}


\begin{keyideabox}
  \textbf{平移口诀：左加右减，上加下减}
  直线 $y=kx+b$ 平移：\\ 1. 左右平移 $h$ 个单位：$x \to (x \pm h)$（左加右减）；\\ 2. 上下平移 $a$ 个单位：整个函数 $\pm a$（上加下减）。
\end{keyideabox}



\begin{question}[points=12]
  将直线 $y = 2x - 1$ 向下平移 $2$ 个单位，得到的直线解析式为
  \begin{choices}
    \item $y = 2x - 3$
    \item $y = 2x + 1$
    \item $y = 2x - 2$
    \item $y = 2(x - 2) - 1$
  \end{choices}
  \paren[A]
  \begin{solutionblock}
    向下平移 2 个单位，即 $b$ 减去 2，故为 $y = 2x - 1 - 2 = 2x - 3$。
  \end{solutionblock}
\end{question}



\begin{question}[points=12]
  将直线 $y = -x + 3$ 向左平移 $1$ 个单位，得到的直线解析式为
  \begin{choices}
    \item $y = -x + 2$
    \item $y = -x + 4$
    \item $y = -(x + 1) + 3$ 即 $y = -x + 2$
    \item $y = -(x - 1) + 3$ 即 $y = -x + 4$
  \end{choices}
  \paren[C]
  \begin{solutionblock}
    向左平移 1 个单位，用 $(x+1)$ 替换 $x$，得 $y = -(x+1) + 3 = -x + 2$。
  \end{solutionblock}
\end{question}



\begin{question}[points=12]
  若一次函数 $y = kx + b$ 的图像与直线 $y = 2x - 3$ 没有交点，且在 $y$ 轴上的截距为 $-9$，则解析式为
  \begin{choices}
    \item $y = 2x - 9$
    \item $y = 2x + 9$
    \item $y = -2x - 9$
    \item $y = 2x - 3$
  \end{choices}
  \paren[A]
  \begin{solutionblock}
    “没有交点”指平行 $\implies k = 2$；在 $y$ 轴上截距为 $-9$ $\implies b = -9$。所以解析式为 $y = 2x - 9$。
  \end{solutionblock}
\end{question}



\begin{question}[points=12]
  将直线 $y = 3x + 2$ 向右平移 $2$ 个单位，得到的直线解析式为
  \begin{choices}
    \item $y = 3x - 4$
    \item $y = 3x + 8$
    \item $y = 3(x - 2) + 2$ 即 $y = 3x - 4$
    \item $y = 3(x + 2) + 2$ 即 $y = 3x + 8$
  \end{choices}
  \paren[C]
  \begin{solutionblock}
    向右平移 2 个单位，用 $(x-2)$ 替换 $x$，得 $y = 3(x-2) + 2 = 3x - 4$。
  \end{solutionblock}
\end{question}


\section*{二、填空题（每题 12 分，共 12 分）}


\begin{question}[points=12]
  将直线 $y = -2x + 5$ 先向上平移 $3$ 个单位，再向右平移 $1$ 个单位，得到的解析式为 \fillin。
  \paren[$y = -2x + 10$]
  \begin{solutionblock}
    先向上平移 3：$y = -2x + 8$。\\ 再向右平移 1：用 $(x-1)$ 替换 $x$，得 $y = -2(x-1) + 8 = -2x + 10$。
  \end{solutionblock}
\end{question}


\section*{三、解答题（共 40 分）}

\needspace{8\baselineskip}

\begin{problem}[points=20]
  已知直线 $y = 2x - 1$。

(1) 写出该直线向下平移 $2$ 个单位后的解析式；
(2) 再将得到的直线向左平移 $1$ 个单位，求平移后直线的解析式，并写出它为何“回到了原解析式”的简要原因。
  \answerarea[50mm]
  \setcounter{solstep}{0}
  \begin{solutionblock}
    \textbf{答案：}(1) $y = 2x - 3$；(2) $y = 2x - 1$，因为沿直线本身的倾斜方向进行了位移\par\medskip
    \solstep{向下平移}{向下平移 2 个单位，截距减小 2。\\ 故平移后的解析式为 $y = 2x - 1 - 2 = 2x - 3$。}
    \solstep{再向左平移}{向左平移 1 个单位，用 $(x+1)$ 替换 $x$。\\ 新解析式为 $y = 2(x+1) - 3 = 2x + 2 - 3 = 2x - 1$。}
    \solstep{简述原因}{可以看出新解析式与原解析式完全一致。\\ 原因是：直线的斜率为 $2$。向左平移 $1$ 个单位相当于向上平移 $2$ 个单位（由斜率定义 $\Delta y / \Delta x = 2$）。\\ 因此向下平移 $2$ 再向左平移 $1$（相当于向上平移 $2$），两个位移刚好沿直线方向抵消，故解析式不变。}
  \end{solutionblock}
\end{problem}


\needspace{8\baselineskip}

\begin{problem}[points=20]
  将直线 $y = -x + 4$ 向上平移 $3$ 个单位，得到新的直线。

(1) 求这新直线的解析式；
(2) 求新直线与坐标轴围成的三角形面积。
  \answerarea[50mm]
  \setcounter{solstep}{0}
  \begin{solutionblock}
    \textbf{答案：}(1) $y = -x + 7$；(2) 面积为 24.5\par\medskip
    \solstep{求新直线解析式}{向上平移 3 个单位，截距增加 3。\\ 故新解析式为 $y = -x + 7$。}
    \solstep{求两轴交点}{令 $y = 0 \implies -x + 7 = 0 \implies x = 7$。与 $x$ 轴交点为 $(7, 0)$，即 $OA = 7$。\\ 令 $x = 0 \implies y = 7$。与 $y$ 轴交点为 $(0, 7)$，即 $OB = 7$。}
    \solstep{求三角形面积}{该三角形是以原点为直角顶点的直角三角形，两直角边长为 $7$。\\ 面积 $S = \frac{1}{2} \times OA \times OB = \frac{1}{2} \times 7 \times 7 = 24.5$。}
  \end{solutionblock}
\end{problem}


\clearpage
\section*{参考答案}


\begin{keyideabox}
  \textbf{选择题}
  1. A  2. C  3. A  4. C
\end{keyideabox}



\begin{keyideabox}
  \textbf{填空题}
  1. $y = -2x + 10$
\end{keyideabox}



\begin{keyideabox}
  \textbf{解答题}
  第 1 题：(1) $y = 2x - 3$；(2) $y = 2x - 1$，因为沿直线本身的倾斜方向进行了位移
第 2 题：(1) $y = -x + 7$；(2) 面积为 24.5
\end{keyideabox}



\end{document}
